\section{두 연산이 만날 때: 환의 출발점}
\label{sec:ring-foundations}

\begin{learninggoal}
이 절을 마치면 다음을 할 수 있다.
\begin{enumerate}[label=(\arabic*)]
  \item 환의 공리를 덧셈군과 곱셈 모노이드의 결합으로 설명한다.
  \item 환, 가환환, 영환, 부분환을 구분한다.
  \item 분배법칙에서 $0a=0$, $(-a)b=-(ab)$ 같은 계산 규칙을 유도한다.
  \item 군과 환의 구조적 차이를 설명한다.
\end{enumerate}
\end{learninggoal}

\subsection{왜 군만으로는 부족한가}

군은 하나의 연산과 그 연산을 되돌리는 방법을 추상화한다. 그러나 정수나 다항식을 다룰 때에는 덧셈만으로 충분하지 않다. 우리는 수를 더할 뿐 아니라 곱하고, 곱셈이 덧셈에 어떻게 분배되는지도 이용한다. 예를 들어
\[
  a(b+c)=ab+ac
\]
는 방정식을 전개하고 인수분해하는 거의 모든 계산의 기반이다.

환은 바로 이 상황을 추상화한다. 하나의 집합 위에 덧셈과 곱셈을 동시에 놓고, 두 연산 사이를 분배법칙으로 연결한다. 이때 덧셈은 항상 아벨군을 이루지만, 곱셈에서는 일반적으로 모든 원소가 역원을 가질 필요가 없다. 이 차이가 환론을 풍부하게 만든다.

\begin{keyidea}
군에서는 모든 원소가 주어진 연산에 대해 가역적이다. 환에서는 덧셈은 항상 가역적이지만 곱셈은 대개 그렇지 않다. 따라서 ``나눌 수 있는 원소'', ``곱하면 사라지는 원소'', ``인수분해되는 원소''를 별도로 연구해야 한다.
\end{keyidea}

\subsection{환의 정의}

이 책에서는 특별히 말하지 않는 한 모든 환은 곱셈 항등원을 가지며, 이후 대부분의 논의에서 곱셈도 가환이라고 가정한다.

\begin{definition}[환]
집합 $R$에 덧셈과 곱셈이라는 두 이항연산이 주어졌다고 하자. 다음 조건을 만족할 때 $R$을 \emph{환}(ring)이라 한다.
\begin{enumerate}[label=(R\arabic*)]
  \item $(R,+)$는 아벨군이다. 덧셈 항등원을 $0$, $a$의 덧셈 역원을 $-a$로 쓴다.
  \item 곱셈은 결합적이며 곱셈 항등원 $1$이 존재한다.
  \item 모든 $a,b,c\in R$에 대해 분배법칙
  \[
    a(b+c)=ab+ac,\qquad (a+b)c=ac+bc
  \]
  이 성립한다.
\end{enumerate}
곱셈까지 가환이면 $R$을 \emph{가환환}(commutative ring)이라 한다.
\end{definition}

\begin{remark}
일부 문헌은 곱셈 항등원을 요구하지 않는 구조도 환이라 부른다. 또한 환 준동형이 $1$을 보존해야 하는지에 대해서도 관례가 다르다. 이 책에서는 환은 항상 $1$을 가지며, 환 준동형은 $1$을 보존한다.
\end{remark}

\begin{example}[기본적인 환]
다음은 모두 가환환이다.
\begin{enumerate}[label=(\alph*)]
  \item 정수환 $\Z$.
  \item 유리수체 $\Q$, 실수체 $\R$, 복소수체 $\C$.
  \item 합동류환 $\Z/n\Z$.
  \item 다항식환 $R[X]$.
  \item 두 가환환 $R,S$의 직접곱 $R\times S$.
\end{enumerate}
행렬환 $M_n(K)$는 $n\ge2$일 때 일반적으로 비가환이다.
\end{example}

\subsection{영환과 공리에서 나오는 계산법}

\begin{proposition}[환의 기본 계산 규칙]
환 $R$의 임의의 $a,b\in R$에 대해 다음이 성립한다.
\begin{enumerate}[label=(\roman*)]
  \item $0a=a0=0$.
  \item $(-a)b=-(ab)=a(-b)$.
  \item $(-a)(-b)=ab$.
  \item 모든 $n\in\Z$에 대해 $(na)b=n(ab)=a(nb)$.
\end{enumerate}
\end{proposition}

\begin{proofidea}
분배법칙을 사용하여 같은 항을 두 방식으로 전개한다. 예를 들어 $(0+0)a$는 한편으로 $0a$, 다른 한편으로 $0a+0a$이다. 양변에서 $0a$를 덧셈으로 소거하면 $0a=0$을 얻는다.
\end{proofidea}

\begin{proof}
분배법칙으로
\[
  0a=(0+0)a=0a+0a
\]
이다. 덧셈군에서 양변에 $-(0a)$를 더하면 $0=0a$이다. $a0=0$도 같다.

또한
\[
  ab+(-a)b=(a-a)b=0b=0
\]
이므로 $(-a)b$는 $ab$의 덧셈 역원이다. 따라서 $(-a)b=-(ab)$이다. 나머지 등식도 같은 방식으로 얻는다.
\end{proof}

$1=0$인 환에서는 위 명제에 의해 임의의 $a\in R$에 대해
\[
  a=a\cdot1=a\cdot0=0
\]
이므로 원소가 하나뿐이다. 이 환을 \emph{영환}(zero ring)이라 한다. 이처럼 $a0=0$을 먼저 공리에서 유도해야 하며, 영환의 결론을 그 성질의 증명에 거꾸로 사용해서는 안 된다.

\begin{pitfall}
환의 곱셈은 군의 연산처럼 모든 원소에 대해 역원을 제공하지 않는다. 따라서 $ab=ac$에서 곧바로 $b=c$라 결론 내리면 안 된다. 예를 들어 $\Z/6\Z$에서
\[
  \overline2\,\overline1=\overline2=\overline2\,\overline4
\]
이지만 $\overline1\ne\overline4$이다.
\end{pitfall}

\subsection{부분환}

\begin{definition}[부분환]
환 $R$의 부분집합 $S$가 $R$의 연산을 그대로 물려받아 환이 되고, $R$과 같은 곱셈 항등원을 가질 때 $S$를 $R$의 \emph{부분환}(subring)이라 한다.
\end{definition}

\begin{proposition}[부분환 판정법]
$R$을 환, $S\subseteq R$라 하자. 다음 조건이 성립하면 $S$는 부분환이다.
\begin{enumerate}[label=(\roman*)]
  \item $1\in S$.
  \item $a,b\in S$이면 $a-b\in S$.
  \item $a,b\in S$이면 $ab\in S$.
\end{enumerate}
\end{proposition}

\begin{proof}
(ii)는 $S$가 덧셈에 대해 부분군임을 보장한다. (i)와 (iii)는 $S$가 곱셈 항등원을 포함하며 곱셈에 대해 닫혀 있음을 보장한다. 결합법칙과 분배법칙은 $R$에서 성립하므로 $S$에도 그대로 제한된다.
\end{proof}

\begin{example}
\[
  \Z\subseteq\Q\subseteq\R\subseteq\C
\]
는 부분환의 사슬이다. 반면 $2\Z\subseteq\Z$는 덧셈과 곱셈에 대해 닫혀 있지만 $1\notin2\Z$이므로 우리의 관례에서는 부분환이 아니다. 그러나 $2\Z$는 이후 정의할 아이디얼의 대표적인 예이다.
\end{example}

\begin{checkpoint}
다음을 스스로 확인하라.
\begin{enumerate}[label=(\alph*)]
  \item $\Z[\sqrt2]=\{a+b\sqrt2:a,b\in\Z\}$는 $\R$의 부분환이다.
  \item $\Q[\sqrt2]=\{a+b\sqrt2:a,b\in\Q\}$는 실제로 체이다.
  \item $\{a+bi:a,b\in\Z\}$는 $\C$의 부분환이다.
\end{enumerate}
\end{checkpoint}
